\documentclass[11pt]{article}
\usepackage[margin=1in]{geometry}
\usepackage{amsmath,amssymb,amsthm,booktabs,array}
\usepackage[T1]{fontenc}
\usepackage{lmodern}
\usepackage[colorlinks=true,linkcolor=blue,citecolor=blue,urlcolor=blue]{hyperref}
\newtheorem{theorem}{Theorem}
\newtheorem{lemma}[theorem]{Lemma}
\newtheorem{corollary}[theorem]{Corollary}
\newcommand{\NN}{\mathbb N_0}
\newcommand{\dist}{\operatorname{dist}}
\newcommand{\diam}{\operatorname{diam}}
\title{Neighborhood two-shadowing for progression subshifts:\newline an extension of Moothathu (2010) answering\newline Meddaugh--Stephens, arXiv:2509.04312v1,\newline Section 5, Question 5.6}
\author{Haobo Ma\\\small ChronoAI Pte Ltd, Singapore
\and Wenlin Zhang\\\small National University of Singapore\\\small ChronoAI Pte Ltd, Singapore}
\date{November 5, 2026}
\begin{document}
\maketitle
\begin{abstract}
Extending Moothathu's progression-exclusion construction (2010),
we answer the open Question 5.6 of Meddaugh and Stephens
(arXiv:2509.04312v1, Section 5, p.~12; 2025) affirmatively by constructing a mixing,
surjective binary forward subshift with neighborhood $2$-shadowing
and a nonempty clopen set containing no nonempty periodic set.
The subshift forbids an all-one arithmetic progression of length
$p^2+3$ and difference $p$, for each positive integer $p$.
Two separated deletion masks from a quadratic-irrational rotation
trace every sufficiently accurate pseudo-orbit. The same construction
is not of quasi-finite type, giving a further affirmative example for
Question 5.3 (Section 5, p.~11), and has no ordinary shadowing.
The same conclusions hold for every progression-length function
$\ell(p)\geq3$ with $\ell(p)/p\to\infty$, without a monotonicity
assumption. Each of the two tracing points can trace at a syndetic
set of times, despite failure of quasi-finite type.
\end{abstract}

\section{Statement and conventions}
Meddaugh and Stephens introduce neighborhood $N$-shadowing in
\cite[Definition 2.1]{MS}. For a compact metric dynamical system
$(X,f)$, this means that for every $\varepsilon>0$ there is $\delta>0$
such that every sequence $(z^i)_{i\in\NN}$ with
$d(f(z^i),z^{i+1})<\delta$ admits a nonempty set $A\subseteq X$ with
$|A|\leq N$, $\diam A<\varepsilon$, and
\[
 \min_{a\in A}d(f^i(a),z^i)<\varepsilon\qquad(i\in\NN).
\]
The choice of $a$ may depend on $i$; the set $A$ does not.
Dense small periodic sets means that every nonempty open set contains
a nonempty compact set $B$ with $f^p(B)=B$ for some $p\geq1$.
Nonemptiness is essential: allowing $B=\varnothing$ would make the
property automatic. Question 5.6 of \cite[Section 5, p.~12]{MS} asks whether weak mixing and neighborhood
$N$-shadowing, for $N>1$, can hold without this property.
For comparison, the finite-cardinality framework of
the author manuscript of \cite[Definition 1.5]{CLMV} uses two-sided pseudo-orbit time and no
separate initial-diameter bound. The definition used here is precisely
\cite[Definition 2.1]{MS}, with natural-time pseudo-orbits and one fixed
tracing set. The affirmative results for mixing sofic subshifts and
systems with specification in \cite[Theorems 3.2, 3.3 and 4.2]{MS}
do not provide the separation requested in Question 5.6.
Our progression constraints retain two-point tracing but exclude
every nonempty periodic subset of the one-cylinder.

The construction has a concrete antecedent: Moothathu
\cite[Proposition 3(iv), p.~131]{Moo} forbids three equally spaced
ones, proves mixing by zero-gap joining, and obstructs repeated
returns to the one-cylinder. A later forward-shift presentation
explicitly credits him \cite[Proposition 4.17, p.~19 of v3]{KLOY}.
The step-dependent forbidden length here permits neighborhood
two-shadowing, which that fixed-length construction does not have.
Irrational-rotation coding of words with bounded constant runs along
each fixed-step progression appears in zeb's 2020 answer on
MathOverflow \cite{MO}; Ronnie Pavlov gave another qualitative
argument there in 2021. Our rotation lemma supplies an explicit
bound suited to two separated deletion masks.
Hereditary deletion and long-zero mixing also occur in
\cite[Example 21 and Proposition 22 of v2]{CKKK}, whose
$\bar d$-shadowing definition controls an upper density of symbol
errors \cite[Section 1.5 of v2]{CKKK}, rather than pointwise
coverage by a fixed nearby pair.

For $p\geq1$ put $L(p)=p^2+3$ and define
\begin{equation}\label{eq:X}
 X=\left\{x\in\{0,1\}^{\NN}:
   \text{for every }s\geq0,\ p\geq1,
   \text{ some }0\leq j<L(p)\text{ has }x_{s+jp}=0\right\}.
\end{equation}
Here and below $\NN=\{0,1,2,\ldots\}$. Give $X$ the metric
$d(x,y)=2^{-\min\{k:x_k\ne y_k\}}$ for $x\ne y$, and $d(x,x)=0$.
Let $\sigma(x)_k=x_{k+1}$.

\begin{theorem}\label{thm:main}
The space $X$ is nonempty and compact, and $\sigma:X\to X$ is
continuous, surjective, and topologically mixing. It has neighborhood
$2$-shadowing. The nonempty clopen cylinder
$U=\{x\in X:x_0=1\}$ contains no nonempty set $B$ satisfying
$\sigma^p(B)\subseteq B$ for any $p\geq1$.
In particular $(X,\sigma)$ answers Question 5.6 affirmatively for every
integer $N\geq2$.
\end{theorem}

The forward-shift convention is explicitly included in
\cite[Section 2]{MS}. The construction is stronger than required:
mixing implies weak mixing, and the obstruction needs neither
compactness of $B$ nor equality of its image with $B$.

\section{Topology, mixing, and the periodic-set obstruction}
Each condition in \eqref{eq:X} depends on finitely many coordinates
and is closed. Thus $X$ is a closed subset of the compact binary
product. It contains the zero point and every point with a single
one. It is shift invariant, and the shift is continuous. Prepending
a zero to a point of $X$ preserves legality: a progression starting
at zero meets that zero, and any other progression translates to
one in the original point. This proves surjectivity.

To prove mixing, fix allowed prefixes $u,v$ of lengths $a,b\geq1$.
For each $n\geq a+\max(a,b)$, place $u$ at $[0,a)$ and $v$ at
$[n,n+b)$, and put zeros everywhere else. Call the resulting point
$x$. A forbidden progression inside either interval would contradict
the legality of the corresponding prefix. A progression meeting both
intervals cannot consist entirely of ones either. If its difference
$p\leq\max(a,b)$, it cannot jump across the zero gap, whose crossing
requires a step at least $n-a+1>\max(a,b)$. If $p>\max(a,b)$, it
contains at most one coordinate in each interval, whereas $L(p)\geq3$.
Consequently $x\in X$. These prefix cylinders form a basis, so every
sufficiently large $n$ maps some point of any given nonempty open
set into any other. This is mixing. Applying the assertion to both
coordinates of two product cylinders at one sufficiently large time
also proves transitivity of $\sigma\times\sigma$, hence weak mixing.

Now suppose $\varnothing\ne B\subseteq U$ and
$\sigma^p(B)\subseteq B$ with $p\geq1$. For $x\in B$, induction gives
$\sigma^{jp}x\in U$, so $x_{jp}=1$ for every $j\geq0$.
The first $L(p)$ of these entries violate \eqref{eq:X}.
This proves the obstruction in Theorem~\ref{thm:main}.

\section{A quantitative rotation lemma}
Write $\{t\}=t-\lfloor t\rfloor$ and
$\|t\|=\min_{m\in\mathbb Z}|t-m|$.

\begin{lemma}\label{lem:rotation}
Let $M\geq1$ be an integer, $\alpha=\sqrt2/M$, and
$P=256M(M+3)$. For every $p\geq1$, $\theta\in\mathbb R$, and
$0\leq a\leq3/4$, some integer $0\leq j\leq Pp$ satisfies
\[
 \{\theta+jp\alpha\}\in[a,a+1/4].
\]
\end{lemma}
\begin{proof}
Put $\kappa=1/(M(M+3))$. We first show
\begin{equation}\label{eq:bad}
 \|k\alpha\|\geq\kappa/k\qquad(k\geq1).
\end{equation}
Choose a nearest integer $m$ to $k\alpha$; it is nonnegative and
$m\leq k\sqrt2/M+1/2$. The integer $2k^2-M^2m^2$ is nonzero,
since $\sqrt2$ is irrational. Therefore
\[
 1\leq |2k^2-M^2m^2|
   =M|k\alpha-m|(k\sqrt2+Mm).
\]
As $2\sqrt2<3$ and $k\geq1$, the last factor is at most
$3k+M/2\leq(M+3)k$, which proves \eqref{eq:bad}.

Consider the $17$ circle points $\{rp\alpha\}$, $0\leq r\leq16$.
One of their $17$ circular gaps has length at most $1/17$.
Subtracting its endpoint indices gives $1\leq q\leq16$ and an
integer $r$ such that
\[
 0<|e|\leq1/17,\qquad e=qp\alpha-r,
 \qquad |e|\geq\kappa/(qp).
\]
The strict inequality follows from irrationality; the lower bound
holds for any integer $r$, by \eqref{eq:bad}.
Set $\tau=a+1/8$ and $h=|e|$. If $e>0$, put
$t=\lfloor\{\tau-\theta\}/h\rfloor$; if $e<0$, put
$t=\lfloor\{\theta-\tau\}/h\rfloor$.
In both cases $t\geq0$, $th<1$, and
$\theta+te-\tau$ is within $h$ of an integer. Hence
$\theta+tqp\alpha$ is within $h<1/8$ of $\tau$ modulo one.
The interval $[a,a+1/4]$ contains that entire circular neighborhood
of $\tau$ and lies in $[0,1]$, so its fractional part belongs to
the desired interval. Finally $j=tq$ satisfies
\[
 j<q/|e|\leq q^2p/\kappa\leq256p/\kappa=Pp.
\]
\end{proof}

\section{Two points trace every pseudo-orbit}
Fix $\varepsilon>0$. Choose an integer $R\geq1$ with
$2^{-(R+1)}<\varepsilon$, and put
\begin{equation}\label{eq:parameters}
 M=32(R+1),\quad \alpha=\frac{\sqrt2}{M},\quad
 P=256M(M+3),\quad D=\max\{R,P^3+2P\},\quad\delta=2^{-D}.
\end{equation}
In particular $R\alpha<1/8$.
For a $\delta$-pseudo-orbit $(z^i)_{i\geq0}$ in $X$, define
the diagonal $c_k=z^k_0$. The pseudo-orbit inequality forces
$(z^i)_{j+1}=(z^{i+1})_j$ for $0\leq j\leq D$.
Iterating this identity gives
\begin{equation}\label{eq:diagonal}
 c_{i+j}=(z^i)_j\qquad(0\leq j\leq D+1).
\end{equation}
For $p\leq P$, each progression of length $L(p)$ in $c$ lies in
a window of span
$(L(p)-1)p=p^3+2p\leq D$. By \eqref{eq:diagonal}, it occurs
in a point of $X$ and therefore contains a zero.

Define two sets of coordinates
\[
 E_0=\{k\geq0:\{k\alpha\}\in[1/8,3/8]\},\qquad
 E_1=\{k\geq0:\{k\alpha\}\in[5/8,7/8]\}.
\]
For $b=0,1$, obtain $a^b$ from $c$ by changing every coordinate
in $E_b$ to zero. For $p\leq P$, deleting ones preserves the zero
already found on every progression. For $p>P$, apply
Lemma~\ref{lem:rotation} with $\theta=s\alpha$ and $a=1/8$ or
$5/8$. The resulting index $j\leq Pp<p^2<L(p)$ puts a zero in
each progression of difference $p$ in $a^b$. Thus $a^0,a^1\in X$.

No coordinate $0\leq k\leq R$ belongs to either mask, since
$0\leq k\alpha\leq R\alpha<1/8$. Consequently $a^0$ and $a^1$
agree on $[0,R]$, so $d(a^0,a^1)\leq2^{-(R+1)}<\varepsilon$.
Moreover each window $[i,i+R]$ avoids at least one mask.
Indeed the circular distance between the two closed quarter intervals
is at least $1/4$, whereas the phase displacement between any two
coordinates of the window is at most $R\alpha<1/8$.
Choose a mask avoided by this window. The corresponding $a^b$ agrees
with $c$ on the whole window, and hence \eqref{eq:diagonal} gives
\[
 d(\sigma^i a^b,z^i)\leq2^{-(R+1)}<\varepsilon.
\]
The fixed set $A=\{a^0,a^1\}$ is nonempty, has at most two points,
has diameter less than $\varepsilon$, and traces the entire
pseudo-orbit. This completes the proof of Theorem~\ref{thm:main}.

\section{Companion consequences and coverage}
\paragraph{Why the fixed-length example does not suffice.}
Let $Y$ be the forward example of \cite[Proposition 4.17]{KLOY}.
It contains singleton configurations but no configuration with
three equally spaced ones. Given $\delta>0$, choose $T\geq2$
with $2^{-(T-1)}<\delta$, and let $z^i$ have its only one at
$(-i)\bmod T$. This is a $\delta$-pseudo-orbit as in the proof
below. If a finite family $a^1,\ldots,a^r$ traced it with error
less than $1$, then for every $j\geq0$ some $a^{c(j)}$ would
have $a^{c(j)}_{jT}=1$. The finite coloring $j\mapsto c(j)$ has
a monochromatic three-term arithmetic progression by van der
Waerden's theorem \cite[Introduction]{KLOY}. One tracer would
therefore contain three equally spaced ones, a contradiction.
Thus this older example fails even finite-family tracing, not
merely the small-diameter requirement.

\begin{corollary}\label{cor:companions}
The system $X$ has exactly one periodic point, the zero point.
It has no ordinary shadowing and is not of quasi-finite type.
\end{corollary}
\begin{proof}
If a periodic point of period $p$ has a one at coordinate $s$,
then it has ones at all $s+jp$, contradicting \eqref{eq:X}.

For failure of ordinary shadowing, fix $\varepsilon=1/2$.
Given $\delta>0$, choose $T\geq2$ with $2^{-(T-1)}<\delta$.
Let $z^i$ have its single one at the least nonnegative residue
of $-i$ modulo $T$. Except when that residue is zero,
$\sigma z^i=z^{i+1}$. At the exceptional transition the shift
is the zero point and the next singleton is at coordinate $T-1$,
so the error is $2^{-(T-1)}$. A single point tracing this pseudo-orbit
within $1/2$ would have $x_i=1$ exactly when $T$ divides $i$.
This violates the constraint for difference $T$.

For the last assertion use \cite[Definition 3.5]{MS}.
Given any proposed quasi-finite-type number $M_0$, choose
$n\geq M_0$ and put $p=n+1$ and $L=L(p)$.
Let $u$ be the word of length $(L-2)p+1$ with ones exactly at
$0,p,\ldots,(L-2)p$, and let $v=1$ and $w=0^n$.
The zero extension of $u$ belongs to $X$: any all-one progression
of at least two terms in it has difference $kp$ for an integer
$k\geq1$, but it has only $L(p)-1<L(kp)$ available ones.
Thus $uw$ and $wv$ are allowed words. For any word $z$ of length
$n$, however, $uzv$ has ones at
$0,p,\ldots,(L-1)p$, since its final coordinate is
$((L-2)p+1)+n=(L-1)p$.
It is forbidden regardless of the entries of $z$.
No $M_0$ can satisfy Definition 3.5.
\end{proof}

\begin{corollary}\label{cor:general}
Let $\ell:\mathbb N_{>0}\to\mathbb N$ satisfy $\ell(p)\geq3$ and
$\ell(p)/p\to\infty$. Replace $L(p)$ in \eqref{eq:X} by $\ell(p)$
and call the resulting forward subshift $X_\ell$. Then
Theorem~\ref{thm:main} and Corollary~\ref{cor:companions} hold for
$X_\ell$. Moreover $X_\ell$ has no mixing number, is not sofic,
and has no specification property in the sense of \cite[Section 2]{MS}.
\end{corollary}
\begin{proof}
The topology, onto property, zero-gap joining, and invariant-subset
obstruction use only $\ell(p)\geq3$, not its growth. For two-point
tracing choose $R,M,\alpha,P$ as in \eqref{eq:parameters}.
Choose $H\geq1$ with $\ell(p)>Pp$ for all $p>H$, and set
\[
 D=\max\left\{R,\max_{1\leq p\leq H}(\ell(p)-1)p\right\}.
\]
The diagonal already satisfies every constraint with $p\leq H$.
For $p>H$, Lemma~\ref{lem:rotation} puts a zero in each mask
by an index $j\leq Pp<\ell(p)$. The window and initial-diameter
arguments are unchanged. The proofs of uniqueness of the periodic
point and failure of ordinary shadowing also apply, using the
constraint of length $\ell(T)$ in the singleton pseudo-orbit argument.

For failure of quasi-finite type, given a proposed threshold $M_0$,
choose $b\geq M_0+1$ and select a positive multiple $q$ of $b$
minimizing $\ell$ among all positive multiples of $b$. Such a minimum
exists because the values are positive integers. Thus
$\ell(kq)\geq\ell(q)$ for every $k\geq1$, even if $\ell$ is not
monotone. Put $n=q-1$, $K=\ell(q)$, and let $u$ have ones exactly
at $0,q,\ldots,(K-2)q$, with length $(K-2)q+1$.
Its zero extension is legal: any all-one progression with at least
two terms has difference $kq$, but $u$ contains only $K-1<\ell(kq)$
ones. With $v=1$ and $w=0^n$, the words $uw,wv$ are legal, whereas
every $uzv$ with $|z|=n$ contains the forbidden $q$-progression
of length $K$. This contradicts Definition 3.5 of \cite{MS}.

A mixing number, as defined after \cite[Proposition 2.3]{MS},
would be a quasi-finite-type number. A mixing sofic subshift has
a mixing number by that proposition, so $X_\ell$ is not sofic.
Finally specification would supply a mixing number: fix
$\varepsilon<1$. Track the symbols of arbitrary allowed words $u,v$
on two orbit segments separated by any gap $n$ at least the
specification constant. A point with the required second segment
exists by prepending zeros to a point beginning with $v$.
Tracking within $\varepsilon$ fixes each symbol, giving a legal
word $uzv$ with $|z|=n$, for every such gap. This is impossible.
\end{proof}

\begin{corollary}\label{cor:syndetic}
The tracing points $a^0,a^1$ in Theorem~\ref{thm:main} and
Corollary~\ref{cor:general} can be chosen so that, for each $b=0,1$,
the set $\{i\geq0:d(\sigma^i a^b,z^i)<\varepsilon\}$ is syndetic:
every interval $[s,s+P]$ of integer times meets it.
\end{corollary}
\begin{proof}
If $\{i\alpha\}\in[1/2,3/4]$, the entire phase window of length
$R$ lies in $[1/2,7/8)$ and avoids $E_0$. If
$\{i\alpha\}\in[0,1/4]$, that window lies in $[0,3/8)$ and
avoids $E_1$. The corresponding point traces at time $i$.
Apply Lemma~\ref{lem:rotation} with $p=1$, starting phase $s\alpha$,
and each of these quarter intervals. In each case a suitable
$i=s+j$ exists with $0\leq j\leq P$.
\end{proof}

\begin{center}
\small
\begin{tabular}{@{}p{0.24\linewidth}p{0.67\linewidth}@{}}
\toprule
Item in \cite{MS} & Coverage of this construction\\
\midrule
Question 5.6 (p.~12) & Answered: yes, with $N=2$, hence every $N\geq2$;
mixing and the stronger invariant-subset obstruction.\\
Question 5.3 (p.~11) & Answered: yes, by the non-quasi-finite-type
argument in Corollary~\ref{cor:companions}.\\
Questions 5.1, 5.2 (p.~11) & Answered: yes at $N=1$, as already
known in \cite[Section 3]{MS}; values $N\geq2$ are not covered.\\
Questions 5.4, 5.5 (p.~12) & The requested general characterizations are
not covered.\\
Fixed-length AP example & Prior proof cited: Moothathu (2010),
with the later forward presentation in \cite{KLOY}. The
step-dependent two-point-tracing extension is distinct.\\
Discussion before Question 5.3 & Not posed; obstruction proved:
syndetic two-point tracing need not imply quasi-finite type
(Corollary~\ref{cor:syndetic}).\\
Converse of Theorem 4.2 & Not asserted in the source; fails for $X$,
which has neighborhood $2$-shadowing but no specification.\\
Proposition 4.3 & Theorem~\ref{thm:main} shows that its periodic-set
hypothesis does not follow from its other hypotheses. It does not
refute the proposition or decide whether that hypothesis is needed
for its mixing conclusion.\\
\bottomrule
\end{tabular}
\end{center}

\paragraph{Proof coverage.}
\begingroup\small\raggedright
All results have written proofs above.
At the pinned \href{https://github.com/the-omega-institute/trureturing/blob/d9bb8a406783c5222849b798f3c2ab738f653d75/D5/S1/Dynamics/Shadowing/ProgressionExample.lean}{\texttt{ProgressionExample.lean}},
the declaration \texttt{quadratic\_\allowbreak progression\_\allowbreak example} states
nonemptiness, compactness, continuity, surjectivity, product
transitivity, neighborhood $2$-shadowing, and the obstruction to
nonempty periodic compact sets. In the same directory,
\href{https://github.com/the-omega-institute/trureturing/blob/d9bb8a406783c5222849b798f3c2ab738f653d75/D5/S1/Dynamics/Shadowing/ProgressionSubshift.lean}{\texttt{ProgressionSubshift.lean}}
contains \texttt{progression\_\allowbreak shift\_\allowbreak mixing} and
\texttt{progression\_\allowbreak no\_\allowbreak invariant\_\allowbreak subset};
\href{https://github.com/the-omega-institute/trureturing/blob/d9bb8a406783c5222849b798f3c2ab738f653d75/D5/S1/Dynamics/Shadowing/ProgressionTwoShadowing.lean}{\texttt{ProgressionTwoShadowing.lean}}
contains \texttt{quadratic\_\allowbreak progression\_\allowbreak neighborhood\_\allowbreak two\_\allowbreak shadowing}.
The same directory's
\href{https://github.com/the-omega-institute/trureturing/blob/d9bb8a406783c5222849b798f3c2ab738f653d75/D5/S1/Dynamics/Shadowing/QuadraticRotationHitting.lean}{\texttt{QuadraticRotationHitting.lean}}
contains \texttt{sqrt\_\allowbreak two\_\allowbreak scaled\_\allowbreak round\_\allowbreak lower}
and \texttt{sqrt\_\allowbreak two\_\allowbreak scaled\_\allowbreak quarter\_\allowbreak hitting},
the lower bound and hitting estimate in Lemma~\ref{lem:rotation}.
\href{https://github.com/the-omega-institute/trureturing/blob/d9bb8a406783c5222849b798f3c2ab738f653d75/D5/S1/Dynamics/Shadowing/ReturnDistanceShadowing.lean}{\texttt{ReturnDistanceShadowing.lean}}
defines \texttt{NeighborhoodTwoShadowing} with one nonempty finite
set, cardinality at most two, small initial diameter, and tracing
at every natural time.
The extra conclusions in Corollaries~\ref{cor:companions}--\ref{cor:syndetic}
have written proofs; in particular the superlinear two-point-tracing
extension and syndetic-time conclusion are distinct from the
quadratic tracing declaration.
The accompanying \texttt{verify.py} checks finite joining,
rotation-mask, and word-obstruction instances with integer arithmetic;
those checks are illustrations, not premises of the infinite arguments.
\par\endgroup

\begin{thebibliography}{6}
\bibitem{CLMV} D. Carrasco-Olivera, K. Lee, C. A. Morales and
H. Villavicencio, \emph{Finite-expansivity and $N$-shadowing},
Bulletin of the Brazilian Mathematical Society, New Series
\textbf{53} (2022), 107--126,
\href{https://doi.org/10.1007/s00574-021-00253-w}{doi:10.1007/s00574-021-00253-w}.
\bibitem{Moo} T. K. Subrahmonian Moothathu,
\emph{Diagonal points having dense orbit},
Colloquium Mathematicum \textbf{120}(1) (2010), 127--138,
\href{https://doi.org/10.4064/cm120-1-9}{doi:10.4064/cm120-1-9}.
\bibitem{KLOY} D. Kwietniak, J. Li, P. Oprocha and X. Ye,
\emph{Multi-recurrence and van der Waerden systems},
Science China Mathematics \textbf{60}(1) (2017), 59--82;
first preprint 7 January 2015,
\href{https://arxiv.org/abs/1501.01491v3}{arXiv:1501.01491v3}
(2 September 2016).
\bibitem{MO} zeb (30 October 2020) and R. Pavlov (28 July 2021),
answers to M. Nielsen, \emph{Binary words that are nonconstant on
long arithmetic progressions}, MathOverflow 375212,
\href{https://mathoverflow.net/questions/375212/binary-words-that-are-nonconstant-on-long-arithmetic-progressions}{question and answers}.
\bibitem{CKKK} M. E. Can, J. Konieczny, M. Kupsa and D. Kwietniak,
\emph{Minimal and proximal examples of $\bar d$-stable and
$\bar d$-approachable shift spaces},
Ergodic Theory and Dynamical Systems \textbf{45}(2) (2025), 396--426;
\href{https://arxiv.org/abs/2308.13967v2}{arXiv:2308.13967v2}
(12 May 2024; first preprint 26 August 2023).
\bibitem{MS} J. Meddaugh and E. Stephens,
\emph{Shadowing With Small Diameter Sets Of Bounded Cardinality},
arXiv:2509.04312v1 (4 September 2025),
\href{https://arxiv.org/abs/2509.04312v1}{arXiv:2509.04312v1}.
\end{thebibliography}
\noindent\small
This note is licensed under Creative Commons Attribution 4.0
International (CC BY 4.0).
\end{document}
